\documentclass[11pt]{article}
\usepackage[T1]{fontenc}
\usepackage[margin=1in]{geometry}
\usepackage{tikz}
\usepackage{pgfplots}
\pgfplotsset{compat=1.18}
\usepackage[utf8]{inputenc}
\usepackage{booktabs}
\usepackage{graphicx}
\usepackage{amsmath}
\usepackage{amssymb}
\usepackage{longtable,tabularx}
\setlength\LTleft{0pt} 
\usepackage{xcolor}
\usepackage{float}
\usepackage{placeins}
\usepackage{ifthen}
\usepackage{natbib}
\usepackage{url}
\usepackage{hyperref}
\graphicspath{{fig/}}  % directory containing figures
\newcommand{\includesvg}[2][]{\includegraphics[#1]{#2}}
\hypersetup{hidelinks}
\setcounter{MaxMatrixCols}{20}

\newboolean{commentsOn}
\setboolean{commentsOn}{true} % or false

\begin{document}

\begin{center}
{\LARGE\bfseries Aircraft System Identification Benchmark Repository Reference}\\[0.75em]
{\large Classical, Grey-Box, Surrogate, and Learning-Based Methods for 6-DOF Aircraft Data}\\[0.75em]
{\normalsize Generated reference document for the \texttt{system\_identification} repository}\\[0.5em]
{\normalsize \today}
\end{center}
\vspace{1em}

\begin{abstract}
This document is the technical reference for the \texttt{system\_identification}
repository. It defines the benchmark motivation, aircraft models, generated
datasets, measurement assumptions, hidden-control cases, implemented
identification methods, code-to-equation traceability, validation metrics, and
current result summaries. The repository focuses on aircraft system
identification for small fixed-wing vehicles when motion capture may be the
only trusted measurement source. The current benchmark suite is a
nonlinear 6-DOF aircraft tier with nominal-plus-unmodeled aerodynamic
terms, motion-capture-derived observations, multiple synthetic maneuver
families, and two real Sport Cub indoor motion-capture flight datasets. The purpose of this reference is reproducibility:
it should be possible to trace each equation, dataset field, result table, and
figure back to repository code and regenerated artifacts.
\end{abstract}

\tableofcontents
\clearpage

\section*{Nomenclature}

\begin{longtable}{@{}l @{\quad=\quad} l@{}}
$N$ & number of discrete samples \\
$n_x, n_u, n_y$ & dimensions of state, input, and output vectors \\
$i$ & discrete-time index ($i = 1, \ldots, N$) \\[3pt]
$\mathbf{x}(t) \in \mathbb{R}^{n_x}$ & continuous-time state vector \\
$\mathbf{u}(t) \in \mathbb{R}^{n_u}$ & continuous-time control input \\
$\mathbf{y}(t) \in \mathbb{R}^{n_y}$ & continuous-time measured output \\
% $\boldsymbol{\theta}$ & vector of unknown model parameters \\
$p_x,p_z$ & inertial longitudinal and vertical position measured by motion capture \\
$V$ & true airspeed (wind axis) \\
$\alpha$ & angle of attack\\
$\gamma$ & flight–path angle \\
$Q$ & body pitch rate about $y$–axis \\
$\theta$ & pitch attitude \\
${\theta_p}$ & vector of constant model parameters \\
$T$ & thrust along $b_x$ \\
$\delta_e$ & elevator deflection \\
$L, D, M$ & lift, drag, and pitching moment \\
$m, J_y$ & mass and pitch moment of inertia \\
$\rho, S$ & air density and reference area \\
$g$ & gravitational acceleration \\
\end{longtable}


\section{Repository Scope and Orientation}

This reference documents the benchmark implemented in the \texttt{system\_identification} repository. It is not a general literature survey and it is not formatted as a journal submission. The purpose is to make the repository auditable: dataset definitions, model equations, method assumptions, generated figures, and result tables should all trace back to code that can be rerun.

The current repository has three main products:
\begin{enumerate}
    \item a reproducible nonlinear 6-DOF aircraft benchmark with stall/residual aerodynamics, synthetic maneuver families, and real indoor motion-capture flight datasets;
    \item a static benchmark website data export and method-plugin interface for future contributed methods.
\end{enumerate}

The primary motivation is small fixed-wing system identification in indoor test environments such as the PURT pylon-racing and fixed-wing UAS testbed setting described by \citet{sribunma2026indoor}. In that setting, the trusted measurement source may be external motion capture rather than onboard air-data or inertial sensors. The benchmark therefore evaluates every method from motion-capture-derived observations: position and attitude are the trusted observables, and any further smoothing or state reconstruction is considered part of the method under test.

\section{Repository Method Inventory}

This section lists only the method families currently represented in repository code or generated benchmark tables. The implementations are centered in \texttt{models/aircraft6dof/comparison\_suite.py}, with the public plugin interface under \texttt{methods/plugins/}. Each method is evaluated by training on its assigned dataset and then performing validation rollouts using the validation initial condition and pilot-command input history. Validation measurements are used for scoring, not for feedback during the rollout, unless a method is explicitly categorized as a measurement-assimilating estimator.

\subsection{Baselines and Fixed-Structure Aircraft Identification}

\paragraph{Nominal grey-box model.}
The nominal grey-box row evaluates the attached-flow 6-DOF aircraft model without the hidden nonlinear aerodynamic residuals. It is a no-fit baseline and appears in the result figures as the reference performance level.

\paragraph{EquationError-LS.}
The equation-error least-squares row fits aerodynamic coefficients from reconstructed state derivatives, the classical regression formulation of aircraft parameter estimation \cite{KleinMorelli2006,morelli_system_2002}. It is fast and interpretable, but it is sensitive to measurement noise because the dependent variables require numerical differentiation.

\paragraph{GreyBoxOEM.}
The grey-box output-error row estimates the physical Sport Cub aerodynamic parameters by optimizing trajectory error through the 6-DOF model with multiple shooting over the training maneuver chunks. It is the main deterministic fixed-structure aircraft identification baseline \cite{jategaonkar2006flight,Morelli2002SIDPAC}, and the same fitted parameters drive the browser free-runs and serve as the UDE base model.

\paragraph{Variational-Mocap.}
The variational/multiple-shooting mocap row adds process-residual flexibility while estimating latent states from mocap outputs. It is included as a practical bridge between fixed-structure output error and process-noise or smoothing formulations.

\paragraph{EKF-ParamID and Fisher-UQ.}
The EKF parameter-identification row estimates parameters with an augmented-state extended Kalman filter \cite{kalman_new_1960,jategaonkar2006flight} and then validates the learned model in open-loop rollout. The Fisher-information row provides local uncertainty and parameter-correlation diagnostics for the fitted fixed-structure model.

\subsection{Linear, Local, Lifted, and Frequency-Domain Predictors}

\paragraph{Linear-SS.}
The linear state-space row fits a discrete affine state-space predictor in the prediction-error tradition \cite{ljung1999system}. It is a strong local predictive baseline, especially near trim or when mocap-derived state reconstruction dominates the error budget, but it should not be interpreted as recovering nonlinear aerodynamic structure.

\paragraph{Model-Stitching.}
The model-stitching row fits local linear models on trim-grid data and blends them using scheduling weights, following the stitched-model and LPV identification practice \cite{Tischler1995SystemIM,toth2010lpv}. It tests whether multiple local models improve over one global linear model across wider initial-condition and maneuver variation.

\paragraph{Subspace-Hankel.}
The subspace row fits a compact linear predictor using Hankel/subspace-style structure. It represents the repository's current state-space realization baseline; it is an ARX-style analogue rather than a full N4SID/MOESP projection method \cite{VanOverscheeDeMoor1996,verhaegen1992subspace,HOKALMAN+1966+545+548}.

\paragraph{Koopman-EDMD.}
The Koopman-EDMD row fits a lifted linear predictor using a fixed nonlinear feature library \cite{koopman1931hamiltonian,williams2015data,schmid2010dmd}. It is included as a predictive surrogate rather than as an aerodynamic-parameter estimator.

\paragraph{Frequency-Welch.}
The frequency-domain row uses Welch/ETFE-inspired spectral estimates \cite{welch1967use}, coherence screening, and a weighted continuous-time linear fit. It is not CIFER \cite{Tischler1995SystemIM} and should be interpreted only as the repository's compact spectral baseline.

\paragraph{Frequency-Stitching.}
The frequency-stitching row applies the same spectral fitting idea locally on trim-grid groups and blends the resulting local continuous-time models. It tests whether the poor global frequency-domain performance is primarily a single-LTI-model limitation.

\subsection{Sparse, Symbolic, Kernel, and Neural Closure Methods}

\paragraph{SINDy.}
The SINDy row \cite{brunton2016discovering,champion2019data} uses a structured coefficient-library formulation so that known low-order aircraft terms can be protected while sparse nonlinear residual terms are selected from a candidate library.

\paragraph{Symbolic-Stepwise.}
The symbolic-stepwise row uses the same candidate-library idea with forward stepwise selection. It is a lightweight symbolic baseline rather than a full genetic-programming symbolic-regression implementation \cite{koza1990genetic,makke_interpretable_2024}.

\paragraph{GP-RBF.}
The GP/RBF coefficient-closure row learns a smooth residual aerodynamic coefficient map with kernel features \cite{rasmussen2006gaussian,hemakumara2013learning}. It represents the repository's current kernel surrogate closure method.

\paragraph{NN-Surrogate.}
The neural coefficient surrogate row learns aerodynamic coefficient residuals from labels available in the synthetic benchmark. Because real flight tests do not provide those coefficient labels directly, this row is an upper-bound diagnostic for what a supervised closure could achieve.

\paragraph{PINN-Closure.}
The PINN-style row \cite{raissi2017physics,lagaris1998artificial} uses flight-dynamics-level residuals, not full CFD or Navier--Stokes constraints, learning aerodynamic coefficient closures consistent with the rollout dynamics.

\paragraph{UDE-Residual and UDE-NN.}
The UDE rows \cite{rackauckas2020universal} retain the known 6-DOF dynamics and learn residual terms targeting the missing aerodynamics. The residual row uses a ridge feature map; the UDE-NN row trains a neural residual through the RK4 rollout (ODE-in-the-loop) on top of the fitted grey-box.

\section{Dataset and Result Artifacts}

The default validation metric in the result tables is the mean range-normalized state RMSE, denoted $J_{\mathrm{val}}$ in this document.

The 6-DOF dataset family splits between local, frequency-excitation, aggressive, and trim-grid cases. The default command
\begin{verbatim}
./results.py simulate-6dof
\end{verbatim}
writes the aggressive nonlinear stall dataset to \texttt{methods/data/aircraft\_6dof\_aggressive/}. The complete 6-DOF dataset family is generated by
\begin{verbatim}
./results.py simulate-6dof --dataset-modes open_loop sine_sweep aggressive trim_grid
\end{verbatim}
where \texttt{open\_loop} is a near-trim small-maneuver case, \texttt{sine\_sweep} is a chirp-like frequency-excitation case, \texttt{aggressive} is the nonlinear stall/recovery stress test, and \texttt{trim\_grid} is a small-deviation grid over speed, angle of attack, and sideslip. The full local 6-DOF workflow is
\begin{verbatim}
./results.py all-6dof
\end{verbatim}
which generates the standard 6-DOF train/validation datasets, runs the implemented 6-DOF baseline methods, exports website data, and refreshes LaTeX-ready artifacts. Training and validation are intentionally separated: local linear, model-stitching, subspace, and frequency-stitching rows train on the trim-grid split; nonlinear residual, surrogate, SINDy, symbolic, and output-error-style rows train on the aggressive split; the nominal row has no fitted training split. Each fitted model is then validated open-loop on the requested validation datasets using only the validation initial condition and pilot-command sequence. The 6-DOF dataset contains inertial position, body velocity, quaternion attitude, body angular rates, pilot commands, actuator-realized commands, noisy direct-state measurements, mocap-style position/quaternion measurements, true aerodynamic coefficients, attached-flow nominal coefficients, and coefficient residual labels.

\section{Legacy 3-DOF Longitudinal Tier}
\label{sec:prob}
Earlier revisions of this repository carried a complete synthetic 3-DOF
longitudinal benchmark (model, datasets, and a twenty-method comparison
suite including CasADi multiple-shooting output error, a filter-error EKF,
and torch ODE-in-the-loop UDE/PINN implementations). Its valuable methods
have been ported to the 6-DOF tier (the filter-error EKF and the
ODE-in-the-loop UDE appear in the method appendix), and the tier itself was
removed to keep the repository focused on the 6-DOF synthetic and
real-flight benchmarks; the full implementation and its results remain
available in the git history. The supplemental method formulations later in
this document originated with that tier and are retained because they are
method theory, independent of the host model.

\section{6-DOF Coupled-Dynamics Benchmark}

The 6-DOF benchmark is the repository's nonlinear coupled-dynamics extension. The state is
\[
\mathbf{x}=
\begin{bmatrix}
x_n & y_e & z_d & u & v & w & q_w & q_x & q_y & q_z & p & q & r
\end{bmatrix}^T,
\]
where $(x_n,y_e,z_d)$ is inertial position, $(u,v,w)$ is body velocity, $(q_w,q_x,q_y,q_z)$ is the unit quaternion mapping body coordinates to inertial coordinates, and $(p,q,r)$ are body angular rates. The input is
\[
\mathbf{u}_{\mathrm{cmd}}=
\begin{bmatrix}
\delta_T & \delta_e & \delta_a & \delta_r
\end{bmatrix}^T.
\]
The simulator applies first-order actuator lag to obtain $\mathbf{u}_{\mathrm{act}}$, but the identification methods are supplied $\mathbf{u}_{\mathrm{cmd}}$ to preserve the practical pilot-command setting.

The rigid-body equations use body-axis velocity and a body-to-inertial quaternion:
\[
\dot{\mathbf{p}}_n = R_{b}^{n}(\mathbf{q})\mathbf{v}_b,\qquad
\dot{\mathbf{v}}_b = \frac{\mathbf{F}_b}{m} + R_{n}^{b}(\mathbf{q})\begin{bmatrix}0&0&g\end{bmatrix}^T - \boldsymbol{\omega}_b\times\mathbf{v}_b,
\]
\[
\dot{\mathbf{q}} = \frac{1}{2}\Omega(\boldsymbol{\omega}_b)\mathbf{q},\qquad
J\dot{\boldsymbol{\omega}}_b = \mathbf{M}_b - \boldsymbol{\omega}_b\times J\boldsymbol{\omega}_b.
\]
The simulator renormalizes the quaternion after every RK4 step.

The truth model computes airdata from body velocity,
\[
V=\|\mathbf{v}_b\|,\qquad \alpha=\tan^{-1}(w/u),\qquad \beta=\sin^{-1}(v/V),
\]
and dynamic pressure $\bar{q}=\frac{1}{2}\rho V^2$. The nondimensional rates used by the aerodynamic model are
\[
\hat{p}=\frac{bp}{2V},\qquad
\hat{q}=\frac{cq}{2V},\qquad
\hat{r}=\frac{br}{2V},
\]
where $b$ is span and $c$ is mean aerodynamic chord. The attached-flow longitudinal coefficients are
\begin{align}
C_{L,\mathrm{att}} &= 0.27 + 5.20\alpha + 0.36\delta_e + 3.10\hat{q}, \label{eq:sixdof_cl_att}\\
C_{D,\mathrm{att}} &= 0.045 + 0.075C_{L,\mathrm{att}}^2 + 0.42\beta^2
    +0.018(\delta_a^2+\delta_r^2)+0.010\delta_e^2, \label{eq:sixdof_cd_att}\\
C_{m,\mathrm{att}} &= 0.030 - 1.05\alpha - 1.15\delta_e - 9.0\hat{q}. \label{eq:sixdof_cm_att}
\end{align}
A smooth stall gate
\[
\sigma(\alpha)=\left(1+\exp\left[-\frac{|\alpha|-\alpha_{\mathrm{stall}}}{\Delta\alpha_{\mathrm{stall}}}\right]\right)^{-1}
\]
uses $\alpha_{\mathrm{stall}}=14^\circ$ and $\Delta\alpha_{\mathrm{stall}}=2.2^\circ$ in the current code. Define
\[
s_\alpha=\operatorname{sign}(\alpha),\qquad
\epsilon_\alpha=\max(|\alpha|-\alpha_{\mathrm{stall}},0),\qquad
C_{L,\lim}=\begin{cases}1.28,&\alpha\ge 0,\\0.95,&\alpha<0,\end{cases}
\]
and
\[
C_{L,\mathrm{post}}=s_\alpha\max(0.28,\;C_{L,\lim}-1.65\epsilon_\alpha),\qquad
\eta_\delta=1-0.58\sigma.
\]
The nonlinear truth lift, drag, and pitching moment are
\begin{align}
C_L &= (1-\sigma)C_{L,\mathrm{att}}+\sigma C_{L,\mathrm{post}}
 +(1-0.45\sigma)\left(0.045\sin(2\alpha)\cos\beta+0.020\delta_e\hat{q}\right), \label{eq:sixdof_cl_truth}\\
C_D &= C_{D,\mathrm{att}}+\sigma\left(0.16+1.45\epsilon_\alpha+5.5\epsilon_\alpha^2\right), \label{eq:sixdof_cd_truth}\\
C_m &= C_{m,\mathrm{att}}
 -\sigma s_\alpha\left(0.18+0.90\epsilon_\alpha\right)
 -0.035\sigma\tanh(4\delta_e). \label{eq:sixdof_cm_truth}
\end{align}
The lateral-directional truth coefficients are
\begin{align}
C_Y &= -0.82\beta + 0.30\delta_r + 0.12\delta_a - 0.35\hat{r}
 +\sigma\left[-0.22\beta|\beta|+0.08\sin(3\alpha)\delta_r\right], \label{eq:sixdof_cy_truth}\\
C_l &= -0.12\beta + 0.42\eta_\delta\delta_a - 0.50\hat{p} + 0.10\hat{r}
 +\sigma(-0.18\beta+0.06\delta_r), \label{eq:sixdof_clroll_truth}\\
C_n &= 0.18\beta - 0.26\eta_\delta\delta_r - 0.08\delta_a - 0.42\hat{r}
 +\sigma(0.10\beta-0.05\delta_a). \label{eq:sixdof_cn_truth}
\end{align}
The nominal 6-DOF model used by \texttt{nominal\_rk4\_step} sets $\sigma=0$ and removes every explicitly nonlinear stall/residual term in \eqref{eq:sixdof_cl_truth}--\eqref{eq:sixdof_cn_truth}. The truth model used by \texttt{rk4\_step} keeps those terms. This nominal/truth split creates the benchmark's hidden residual-aerodynamics identification problem.

Body aerodynamic forces are then
\[
\begin{bmatrix}F_X&F_Y&F_Z\end{bmatrix}^T
= \bar{q}S\begin{bmatrix}C_X&C_Y&C_Z\end{bmatrix}^T,
\]
with
\[
C_X=-C_D\cos\alpha+C_L\sin\alpha,\qquad
C_Z=-C_D\sin\alpha-C_L\cos\alpha.
\]
Moments are
\[
\mathbf{M}_{\mathrm{aero}}=\bar{q}S
\begin{bmatrix}
bC_l & cC_m & bC_n
\end{bmatrix}^T.
\]
The propeller model is
\[
T=T_{\max}\delta_T^{1.45},\qquad
\mathbf{F}_{\mathrm{prop}}=\begin{bmatrix}T&0&0\end{bmatrix}^T,\qquad
\mathbf{M}_{\mathrm{prop}}=\begin{bmatrix}0&\ell_TT&0\end{bmatrix}^T,
\]
with $T_{\max}=8.8\,\mathrm{N}$ and $\ell_T=0.035\,\mathrm{m}$. The coded body loads are
\[
\mathbf{F}_b=(1+0.22\delta_T)\mathbf{F}_{\mathrm{aero}}+\mathbf{F}_{\mathrm{prop}},
\qquad
\mathbf{M}_b=(1+0.22\delta_T)\mathbf{M}_{\mathrm{aero}}+\mathbf{M}_{\mathrm{prop}}.
\]
The dataset stores both the truth coefficients and an attached-flow nominal coefficient set; their difference is the hidden nonlinear aerodynamic residual. Four maneuver families are available in the generator: a near-trim \texttt{open\_loop} small-maneuver case, a \texttt{sine\_sweep} frequency-excitation case, an \texttt{aggressive} nonlinear stall/recovery case, and a \texttt{trim\_grid} case that samples local small-deviation responses across speed, angle of attack, and sideslip. The 6-DOF comparison assigns every method an explicit training split and then validates the trained model across the maneuver family. The local-model methods train on \texttt{trim\_grid} because they require repeated small perturbations near operating points; global nonlinear residual and surrogate methods train on \texttt{aggressive}; the nominal baseline uses no fit. The 6-DOF heatmaps report all validation datasets as separate columns, with the leftmost column giving each method's mean score across the available 6-DOF benchmark conditions.

A single observation policy is evaluated: the method receives motion-capture position and quaternion attitude. Body velocity is reconstructed by differentiating mocap position and rotating the result into the body frame, while angular rate is reconstructed by differentiating the quaternion and solving the local quaternion-rate equation. These reconstructed velocity and rate channels are intentionally noisier than an idealized full-state sensor; any additional smoothing or state estimation a method applies is part of that method. Earlier revisions also reported an idealized direct-state diagnostic, but on the real datasets both channels derive from the same motion-capture sensor, so the distinction was retired.

\subsection{Real Flight Records: Lead-In Segmentation and Trim-Offset Calibration}
\label{subsec:real_segmentation}

In addition to the synthetic maneuver families, the 6-DOF benchmark runs the
same baseline methods on real Sport Cub S2 motion-capture flight records
collected in the PURT indoor testbed \citep{sribunma2026indoor}. Two
collections are included. The 2026-04-17 collection consists of short piloted
maneuvers (elevator and aileron 3-2-1-1 sequences, doublets, and throttle
steps) tracked at 100\,Hz and packaged as train and validation splits of
maneuver segments. The 2026-05-22 collection consists of complete piloted
laps recorded as ROS~2 bags: the aircraft flies a stabilized circuit with the
SAFE inner loop engaged, the pilot disengages the inner loop mid-facility,
performs a maneuver on the bare airframe, and re-engages stabilization. The
transmitter mode channel is recorded alongside the stick channels, so the
controller state is known at every sample. A converter resamples the mocap
odometry and 53\,Hz stick stream to a uniform 240\,Hz grid (the native
Qualisys odometry rate; one autonomous bag streams 100\,Hz) (the stick-to-PWM
calibration was verified exact, $|r|>0.999$, against the serial echo on every
manually flown bag) and an automatic segmenter classifies every sample as
ground, stabilized, or manual: ground samples are detected from motion-capture
altitude with hysteresis and excluded; stabilized samples are excluded from the
bare-airframe splits because the hidden inner loop reshapes the surface
commands; quality-gated manual windows become the identification segments,
assigned round-robin across flights so one model is learned from all runs.
Bare-airframe trim is estimated per flight by pooling quasi-steady manual
samples (low angular rate, low stick rate), because each run can carry a
different transmitter trim. Real records differ from the synthetic datasets
in two ways that matter for open-loop scoring. First, each segment begins with
an input-free dwell of varying length before the pilot applies the maneuver
input. Because the airframe is open-loop unstable, the rollout error
accumulated during this dwell grows at the rate of the unstable mode
regardless of model quality, so the dwell measures initial-condition
sensitivity rather than identification accuracy. Second, the transmitter trim
is not exactly zero, so the recorded stick commands carry unknown constant
offsets relative to the model's neutral control surfaces; left uncorrected,
these offsets bias every fitted control derivative.

The benchmark therefore decomposes each real record into a \emph{lead-in}
segment and a \emph{control-actuation} segment, following standard flight-test
practice in which maneuver records are windowed around the excitation and
constant bias parameters are estimated as part of data compatibility analysis
\citep{KleinMorelli2006, jategaonkar2006flight, Tischler1995SystemIM}. The
lead-in length is detected automatically per segment as the first sample at
which any command channel departs from its initial level by more than 15\% of
that channel's range within the record (with an absolute floor to reject
noise-only channels). The lead-in is used only for trim calibration; the
control-actuation segment is used for dynamics, providing the fitting data on
the training split and the open-loop scoring window on the validation split.

The lead-in calibrates two quantities. The validation initial condition is
estimated by a local linear regression over the lead-in evaluated at the
actuation onset, rather than read from a single noisy measurement frame.
Constant input offsets $\mathbf{b}$ are estimated by reconciling the nominal
model with the quasi-steady lead-in: with $\bar{\mathbf{x}}$ and
$\bar{\mathbf{u}}$ the lead-in mean state and command and
$\dot{\mathbf{x}}_{\mathrm{obs}}$ the secant state derivative across the
lead-in, the offsets solve the regularized least-squares problem
\begin{equation}
    \frac{\partial \mathbf{f}}{\partial \mathbf{u}}
    \Big|_{\bar{\mathbf{x}},\bar{\mathbf{u}}}\,\mathbf{b}
    \approx \mathbf{f}(\bar{\mathbf{x}},\bar{\mathbf{u}})
    - \dot{\mathbf{x}}_{\mathrm{obs}},
    \label{eq:input_bias}
\end{equation}
and the corrected commands $\mathbf{u}-\mathbf{b}$ are supplied to all
methods. The kinematic rows of $\mathbf{f}$ do not depend on the inputs and
drop out of the normal equations. Because the transmitter trim is a single
physical setting per flight, the normal equations are pooled per flight
(across all maneuver segments recorded in that flight), weighted by lead-in
length, rather than solved per segment; per-segment estimates were observed to
absorb local model error into the offsets. The estimate is clipped to $\pm0.2$ command units, and the usual
caveat of data compatibility analysis applies: the recovered offsets are
\emph{effective} biases that absorb any steady model error (for example,
thrust-map mismatch appearing as a throttle offset) along with the physical
transmitter trim \citep{KleinMorelli2006}. Within the benchmark structure the
fitted-model rows are unaffected by this confounding because their regressors
include constant columns; the correction primarily re-references the
physics-based rows to the aircraft's actual trim. Because real maneuver
records have very different lengths, the actuation segments are
rectangularized by chunking rather than truncation: a common window length is
chosen from the shortest usable record (clamped to a benchmark range of
0.6--1.2\,s) and longer maneuvers contribute several consecutive windows, with
follow-on window initial conditions estimated from a local fit over the
trailing samples of the preceding window. No flight data is discarded to the
shortest segment. The segmentation and calibration are applied identically to
the training and validation splits, each split using only its own lead-in
data, and can be disabled with \texttt{--no-onset-trim} and
\texttt{--no-input-bias} for ablation.

\paragraph{Ground-roll data and rolling-resistance identification.}
The segmenter also exports the rolling ground data instead of discarding it,
classified into takeoff rolls (throttle advance to the instant of liftoff),
landing rollouts, and taxi, plus low passes below approximately one wingspan
tagged as a (small) ground-effect set. Ground data is highly susceptible to
floor disturbances — wheel debris, surface seams, handling contact, and the
taildragger gear's tendency to ground-loop — so every
sample carries a robust outlier flag before any fitting. Yaw-rate bursts above
1.5\,rad/s on the ground are flagged as spin-outs, whose wheel-scrub forces
invalidate the straight-rolling friction model (deliberate taxi turns are
flagged by the same test, and are equally invalid for that model). Detectors based on
differentiated quantities fail here: acceleration z-scores flag the legitimate
acceleration ramp of a takeoff roll, and jerk z-scores flag the resampling
knots of segments recorded at lower odometry rates. The detector instead
flags departures of ground speed from a short moving-average fit, using a
median/MAD modified z-score \emph{with an absolute magnitude floor}
(0.15\,m/s): rolling contact has a small-amplitude periodic texture from the
wheels and floor seams that a pure z-score flags wholesale once the robust
scale collapses, while genuine debris strikes also clear the absolute floor.
An attitude gate flags departures from the per-window median attitude (the
taildragger rests nose-up, so absolute pitch is not an anomaly), and flags are
dilated to cover transient shoulders; about 4\% of rolling samples are
rejected.

The ground-roll model is $\dot{V} = k_T T(\delta_T)/m - \mu_{\mathrm{eff}} g -
c_v V^2 - k_{\mathrm{batt}}\,\tau\,T(\delta_T)/m$ with $T(\delta_T) = T_{\max}
\delta_T^{1.45}$ the nominal thrust map and $\tau$ the session time. In this
session the pilot held essentially one throttle level on the ground, so $k_T$
and $\mu_{\mathrm{eff}}$ are not separately identifiable from ground data
alone; the analysis adopts the nominal map as a stated prior ($k_T = 1$) and
reports the sensitivity $\partial\mu_{\mathrm{eff}}/\partial k_T$ so the
assumption is transparent. The $k_{\mathrm{batt}}$ term models battery sag:
the airframe flies on a 1S LiPo whose prop power visibly decays over the
session, and because the ground windows span the session at a common throttle
level, session time provides the thrust contrast that throttle does not.
A second battery stage separates the effect per flight, and is the physically
meaningful one: pack endurance is only 5--10 minutes, so batteries are swapped
or recharged between flights and the session-wide term is at best a crude
trend across packs rather than one battery discharging. Each flight starts at
its own pack voltage and the voltage droops within the flight, lowering the
achievable motor rpm. Voltage is not telemetered, so with
$\mu_{\mathrm{eff}}$ and $c_v$ anchored by the pooled fit, the thrust-side
residual is regressed per flight on $[T/m,\; t_{\mathrm{flight}}\,T/m]$,
giving an initial thrust scale $k_f$ (an initial-voltage proxy) and a
within-flight decay rate for every flight with sufficient ground samples.
These per-flight estimates carry two stated caveats: flights whose ground
data is takeoff-only leave $t\,T/m$ nearly collinear with $T/m$ at the
session's single throttle level, and landing-rollout samples mix propeller
windmilling drag into the apparent late-flight thrust loss, so the per-flight
numbers are meaningful primarily where both takeoff and landing rolls exist. A
Huber iteratively-reweighted fit over the outlier-screened, yaw-rate-gated
rolling samples yields $\mu_{\mathrm{eff}} \approx 0.10$ (plausible for small
rubber wheels on a gymnasium floor; sensitivity $0.16$ per unit $k_T$),
$c_v \approx 0.005\,\mathrm{m}^{-1}$, and a thrust loss of roughly 7\% per
30 minutes of session time. Figure~\ref{fig:ground_roll} shows the longest
takeoff roll with its outlier flags and the pooled fit.

\begin{figure}[htbp]
    \centering
    \includesvg[width=0.9\linewidth]{ground_roll_takeoff}
    \caption{Automatically segmented takeoff roll from the 2026-05-22 Sport
    Cub session. Top: ground speed and scaled throttle with robust outlier
    flags (red). Bottom: measured longitudinal acceleration against the pooled
    Huber ground-roll fit with rolling resistance, quadratic drag, and battery
    sag under the nominal-thrust prior.}
    \label{fig:ground_roll}
\end{figure}

\paragraph{Hidden SAFE inner-loop identification.}
Free-run prediction through stabilized flight requires a model of the hidden
Horizon Hobby SAFE inner loop, because in stabilized mode the pilot is not
commanding surface deflections at all: the stick commands attitude (and
throttle passes through), with envelope protection limiting the achievable
attitude and the usual surface saturation. Per attitude axis the controller is
modeled as
$\delta = K_p\left(\mathrm{sat}(c\,\delta_{\mathrm{stick}},\pm\ell) -
\vartheta\right) - K_d\,\omega + b$,
where $c$ is the stick-to-attitude-command scale, the clip at $\pm\ell$ is the
envelope protection, $\vartheta$ and $\omega$ are the axis attitude and rate,
and the output saturates at the surface throws; the rudder is modeled as
passthrough plus yaw damping. The model is nonlinear only through $(c,\ell)$,
so identification is staged: a coarse grid over $(c,\ell)$ wraps a
Huber-robust linear fit of $(K_p, K_d, b)$ given the saturated attitude
command. The effective surfaces are not measured, but inverse dynamics
recovers them sample by sample: the observed angular accelerations give the
aerodynamic moments, the airframe moment model is affine in each surface, and
inverting that relation yields the deflection the airframe must have seen —
the 6-DOF analogue of the benchmark's hidden-controller rows. The current
estimate is provisional: the inversion is conditioned on the nominal
airframe's surface effectiveness, the double-differentiated motion-capture
attitude is noisy, and the recorded laps barely probe the envelope limits, so
the gains improve as the identified airframe replaces the nominal model and as
envelope-probing stabilized data is collected.

\paragraph{Segment-anchored scoring and free-run prediction.}
The full-flight explorer view scores each airborne segment with a rollout
whose initial condition is estimated by a local fit at that segment's start —
the same convention used during training, so identification cost functions
never begin from an already-diverged state on an open-loop unstable airframe.
From any selected segment a free run continues to the next ground contact
with no further state resets, handing off between two models on the recorded
flight-mode channel with the state carried continuously across each switch:
while the pilot flies manual the bare airframe integrates the recorded
sticks re-referenced by the per-flight trim bias, and while SAFE is engaged
the free run uses a \emph{directly identified closed-loop model} (the
\emph{direct method} of closed-loop identification \cite{ljung1999system})
ridge-fitted on the tracked stabilized data rather than composing the bare
airframe with the inverse-dynamics controller estimate: the closed loop is
abundantly excited by the lap data and is exactly what stabilized segments
replay. The stabilized spans receive the same train/validation discipline
as the manual maneuver windows: they are cut into ten-second windows
assigned round-robin per flight with every third window held out, the model
fits on the train windows only, and the held-out windows score it by
free-run position error,
whereas the decomposition route inherits the weakly conditioned controller
estimate (on five-second stabilized free-runs the direct fit is roughly
four-fold better). The regression is restricted to the heading- and
position-invariant state --- body velocities, roll, pitch, and body rates,
with the heading \emph{increment} as an additional target --- because flight
dynamics do not depend on position or heading, and a global affine map
fitted on raw position/quaternion states cannot represent the
heading-dependent position kinematics: its free runs wander off the
operating point within seconds. Heading then integrates the fitted
increment and position integrates the rotated body velocity exactly, which
lets stabilized free runs fly whole laps; on ten-second held-out stabilized
windows this symmetry-respecting fit cuts the five-second position error
roughly three-fold relative to the raw-state affine fit (about 3\,m versus
8\,m). The
explicit controller identification of the previous paragraph is retained for
interpretability and for eventual composition with identified airframes. Model rollouts are computed
on the fly in the browser from exported model parameters (the attached-flow
nominal model, affine state-space and residual weights, and controller
gains), which avoids exporting a trace per segment--method pair and allows
the initial condition to be placed anywhere in the flight; small
NumPy-trained networks deploy the same way as JSON weight specifications
evaluated by a generic forward pass, while torch-trained models deploy as
ONNX graphs evaluated in-browser, with the same integration loop in either
case.

\paragraph{Physical grey-box output-error fit.}
The one-step affine predictors that dominate the synthetic leaderboards
extrapolate poorly through aggressive manual segments: their free runs are
anchored to the regression's operating envelope rather than to flight
physics. The benchmark therefore also fits the repository's lumped-parameter
Sport Cub grey-box (twelve Euler states; lift, drag, and side-force
coefficient expansions plus moment-acceleration coefficients with
$b/2V$ and $\bar c/2V$ rate terms) by segment-wise output error
\cite{jategaonkar2006flight}: every
bias-corrected manual training chunk is rolled out with a fixed-step RK4
integrator from its measured initial state, and the 22 aerodynamic
parameters jointly minimize position and attitude residuals (weighted by the
spec's output sigmas) across all chunks, constrained to physical bounds by a
trust-region reflective solver warm-started at handbook initials. On the
2026-05-22 manual windows the fitted model roughly halves the held-out
rollout NRMSE relative to the bound-clipped handbook initial guess
(0.14 versus 0.25), with lift-curve and thrust parameters pushing against
their upper bounds --- consistent with the ground-roll finding that the
nominal thrust constant underestimates the installed power. Because the
model is a compact closed-form ODE, the same fitted parameters export
directly to the browser, where the explorer integrates it with the identical
RK4 step for segment scores and free-run predictions. Known structural
limitations, in roughly decreasing order of expected impact for this
airframe: no actuator/servo dynamics (a 50--100\,ms lag is identifiable at
240\,Hz and its absence biases control derivatives low), a static thrust map
with no airspeed or battery-stage dependence even though the ground-roll
analysis quantifies a per-flight thrust stage, no propeller-slipstream
augmentation of tail effectiveness, and no $C_{L_q}$ or
$\dot\alpha$ terms. A ground test resolves whether gyro
augmentation contaminates the ``manual'' (Experienced-mode) segments:
rotating the airframe by hand in Experienced mode produces no surface
response, so no rate feedback is active there and the fitted rate-damping
coefficients are genuine bare-airframe aerodynamic derivatives.

\paragraph{Identifiability and uncertainty.} Cram\'er--Rao lower bounds
from the output-error Jacobian at the solution (residual coloring
uncorrected, so the bounds are optimistic) split the 22 parameters into
three classes. Well identified (1--8\% relative standard deviation): the
lift line ($C_{L_0}$ 1.0\%, $C_{L_\alpha}$ 3.1\%), all primary damping
and stiffness terms ($K_{M_q}$ 2.7\%, $K_{L_p}$ 3.5\%, $K_{N_\beta}$
2.5\%), and the control effectivenesses. Weakly identified: $K_T$
(14.8\%), $K_{L_{\delta a}}$ (17.2\%), $K_{N_0}$ (27\%). Unidentifiable:
$C_{D_{CL^2}}$ at 368\% --- the formal confirmation of its physical
zero-bound pin. Four parameter pairs are strongly coupled
($|r|>0.9$): $C_{L_0}\!\leftrightarrow\!C_{L_\alpha}$ and
$K_{M_0}\!\leftrightarrow\!K_{M_\alpha}$ (the laps occupy a narrow
angle-of-attack band, so intercept and slope trade off), and
$K_{L_0}\!\leftrightarrow\!K_{L_{\delta r}}$,
$K_{N_0}\!\leftrightarrow\!K_{N_{\delta r}}$ (the rudder rides near a
constant trim, aliasing its effectiveness with the moment biases). The
remedies are input-design items: wider angle-of-attack excitation and
rudder doublets about varied trims. The full covariance diagnostics are
written into the fitted-parameter artifact.

\paragraph{Actuator lag and battery decay: tested, not identifiable here.}
A 25-parameter extension adding first-order servo and motor lag states and
a linear battery thrust-decay term was fitted and gated on held-out data:
all three new parameters pinned at their bounds (both lags at their minima,
the decay at its maximum) and validation NRMSE \emph{worsened} from 0.106
to 0.130. The mechanism is observability: the 3-2-1-1 excitation at
2--4\,rad/s cannot see a $\sim$30\,ms servo pole, and within-flight
maneuver placement confounds the per-flight-time decay term, which then
absorbs unrelated trends. The lag-augmented model is retained in the code
for reuse once dedicated excitation exists; the battery effect remains
quantified independently by the ground-roll analysis, and folding those
throttle-step ground runs into the flight fit is the identified path to a
recoverable $K_B$.

\subsection{6-DOF Baseline Methods}

\paragraph{6DOF-NominalGreyBox.}
The 6-DOF nominal grey-box row propagates the repository's attached-flow quaternion rigid-body model with the supplied pilot-command history. It has no fitted parameters and is used as the no-fit baseline against the hidden nonlinear stall/residual aerodynamic truth model. The row appears only for synthetic scenarios: on real flights the truth-minus-residual interpretation does not exist and the attached-flow model is simply a wrong model, so it is omitted there (it remains the internal base model of the residual methods).

\paragraph{6DOF-LinearSS.}
The 6-DOF linear state-space row fits an affine one-step predictor,
\[
\mathbf{x}_{k+1}=A\mathbf{x}_k+B\mathbf{u}_{\mathrm{cmd},k}+\mathbf{c},
\]
using ridge-regularized least squares. During validation the model is rolled out open-loop from the validation initial condition and pilot-command sequence.

\paragraph{6DOF-RidgeResidual.}
The 6-DOF residual row first propagates the nominal RK4 model and then adds a ridge-fitted one-step residual correction,
\[
\mathbf{x}_{k+1}=\Phi_{\mathrm{nom}}(\mathbf{x}_k,\mathbf{u}_{\mathrm{cmd},k})
+W^T[\mathbf{x}_k^T,\mathbf{u}_{\mathrm{cmd},k}^T,1]^T.
\]
This tests the same grey-box idea used in UDE-style methods, but with a linear residual so that the 6-DOF benchmark remains fast and deterministic. The residual absorbs actuator lag and the hidden difference between the attached-flow nominal model and the stalled nonlinear truth model.

\IfFileExists{fig/generated_aircraft6dof_validation_score_comparison.svg}{%
\begin{figure}[hbt!]
    \centering
    \includesvg[width=0.95\linewidth]{generated_aircraft6dof_validation_score_comparison}
    \caption{6-DOF baseline method validation scores (full-state mean NRMSE over open-loop rollouts). Lower is better.}
    \label{fig:aircraft6dof_validation_score_comparison}
\end{figure}
}{}

\IfFileExists{fig/generated_aircraft6dof_train_time_accuracy_tradeoff.svg}{%
\begin{figure}[hbt!]
    \centering
    \includesvg[width=0.95\linewidth]{generated_aircraft6dof_train_time_accuracy_tradeoff}
    \caption{Training-time versus validation-error tradeoff for the 6-DOF benchmark. Methods with mean-state NRMSE above 1.0 are treated as failed rollouts and summarized in the callout so the useful region remains visible. The red dashed horizontal line gives the attached-flow nominal baseline, which has no fitted training solve. Marker area increases with validation rollout time.}
    \label{fig:aircraft6dof_train_time_accuracy_tradeoff}
\end{figure}
}{}

\IfFileExists{fig/generated_aircraft6dof_method_score_heatmap.svg}{%
\begin{figure}[hbt!]
    \centering
    \includesvg[width=0.95\linewidth]{generated_aircraft6dof_method_score_heatmap}
    \caption{6-DOF benchmark heatmap sorted by mean validation score. The score is full-state mean NRMSE; methods initialize and train from position/attitude-derived state estimates. Columns report the mean across 6-DOF validation datasets followed by the open-loop, sine-sweep, aggressive nonlinear stall, and trim-grid validation datasets. The training split is method-specific and reported in Table~\ref{tab:aircraft6dof_method_comparison}; local-model methods train on the trim grid, global nonlinear residual/surrogate methods train on the aggressive split, and the nominal baseline has no fitted training solve.}
    \label{fig:aircraft6dof_method_score_heatmap}
\end{figure}
}{}

\IfFileExists{fig/generated_aircraft6dof_validation_trajectory_overlay.svg}{%
\begin{figure}[hbt!]
    \centering
    \includesvg[width=0.95\linewidth]{generated_aircraft6dof_validation_trajectory_overlay}
    \caption{Representative 6-DOF validation rollout for the best implemented state-predictor baselines. Rollouts use only the validation initial condition and pilot-command sequence after training; the panels include position, speed, angle of attack, stall-gate activation, and body rates.}
    \label{fig:aircraft6dof_validation_trajectory_overlay}
\end{figure}
}{}

Website-ready benchmark data are exported by
\begin{verbatim}
./results.py web-data
\end{verbatim}
into \texttt{site/public/data/}. The static site can be served locally with
\begin{verbatim}
./results.py serve-site
\end{verbatim}
The method-contribution workflow is intentionally two-phase: public pull requests add method code and pass plugin smoke checks, while full benchmark results are regenerated later in a trusted maintainer commit.

\section{Supplemental Method Formulations}
This section keeps equation-level details for implemented method families whose
benchmark results are generated by the repository. Legacy standalone figures
and tables that are not regenerated by the current pipeline have been removed.

\subsection{Output Error Method}
This formulation belongs to the class of prediction error methods (PEM), where the objective is to minimize the discrepancy between model-predicted outputs and measured data. The Output Error (OE) method is a special case that assumes deterministic system dynamics, with all uncertainty entering through measurement noise.

Starting from the nonlinear system,
\begin{align}
\dot{x}(t) &= f\big(x(t),u(t);\theta_p\big),\\
y(t) &= h\big(x(t);\theta_p\big) + v(t),
\end{align}
the predicted output $\hat{y}(t,\theta_p)$ is obtained by numerically integrating the state equations for a given parameter vector $\theta_p$ (see Section~\ref{sec:prob}).

Assuming additive Gaussian measurement noise,
\begin{align}
y_k = \hat{y}_k(\theta_p) + \varepsilon_k, \qquad \varepsilon_k \sim \mathcal{N}(0, R),
\end{align}
the maximum likelihood estimation (MLE) problem becomes
\begin{align}
J_{\text{MLE}}(\theta_p, R) =
\frac{1}{2} \sum_{k=1}^{N} 
\left(y_k - \hat{y}_k(\theta_p)\right)^\top R^{-1} 
\left(y_k - \hat{y}_k(\theta_p)\right)
+ \frac{N}{2} \log \det R.
\end{align}
In this work, $R$ is assumed diagonal and estimated jointly with $\theta_p$ via a log-parameterization, resulting in adaptive weighting across output channels.

The resulting nonlinear program is constructed symbolically using CasADi and solved using IPOPT, with state trajectories generated through a fixed-step fourth-order Runge--Kutta (RK4) integrator embedded within the optimization.

Two OE formulations are considered: single shooting (SS) and multiple shooting (MS).

\subsubsection{Single--Shooting Formulation}

In the single--shooting implementation, the full state trajectory is generated by
integrating the nominal longitudinal model over the complete experiment horizon.
The aerodynamic parameter vector, the initial state, and the measurement-noise
standard deviations are estimated jointly. The decision vector is
\[
    w_{\mathrm{SS}}
    =
    \begin{bmatrix}
        \boldsymbol{\theta}_p^\top &
        \mathbf{x}_0^\top &
        \log\boldsymbol{\sigma}^\top
    \end{bmatrix}^{\!\top}.
\]
Starting from the optimized initial condition $\mathbf{x}_0$, the trajectory is
propagated using a fixed-step fourth-order Runge--Kutta map,
\[
    \hat{\mathbf{x}}_{k+1}
    =
    \Phi(\hat{\mathbf{x}}_k,\mathbf{u}_k;\boldsymbol{\theta}_p,\Delta t),
    \qquad k=0,\ldots,N-1.
\]
In the oracle state-measurement comparison, the predicted output is
taken as $\hat{\mathbf{y}}_k=\hat{\mathbf{x}}_k$. In the motion-capture comparison,
the same latent trajectory is instead passed through the measurement model
\[
\hat{\mathbf{y}}_{\mathrm{mocap},k}
=h(\hat{\mathbf{x}}_k)
=
\begin{bmatrix}
\hat{p}_{x,k} & \hat{p}_{z,k} & \hat{\alpha}_k+\hat{\gamma}_k
\end{bmatrix}^{T},
\]
where $\hat{p}_x$ and $\hat{p}_z$ are obtained by integrating
$\dot{p}_x=\hat{V}\cos\hat{\gamma}$ and
$\dot{p}_z=\hat{V}\sin\hat{\gamma}$. This output-error formulation avoids
using numerically differentiated mocap data as the dependent variable. The negative log-likelihood
objective, assuming independent Gaussian noise in each channel, is
\[
    J_{\mathrm{SS}}
    =
    \sum_{j=1}^{n_y} N\log\sigma_j
    +
    \frac{1}{2}
    \sum_{k=0}^{N-1}
    \sum_{j=1}^{n_y}
    \frac{\left(\hat{y}_{k,j}-y_{k,j}\right)^2}{\sigma_j^2}.
\]
Thus, the single--shooting OEM problem is
\[
    \min_{\boldsymbol{\theta}_p,\,\mathbf{x}_0,\,\log\boldsymbol{\sigma}}
    J_{\mathrm{SS}}.
\]
While this formulation is compact, the entire predicted trajectory relies on a single integration from the initial state, rendering the optimization sensitive to parameter initialization and prone to error accumulation over long horizons. Consequently, the Equation Error Method is often employed as a precursor to provide a robust initial guess for the OEM. The single-shooting output error approach remains standard practice for low-turbulence flight tests and is implemented in widely used toolkits such as SIDPAC \cite{Morelli2002SIDPAC}; when process noise (turbulence) is significant, the filter-error method, which embeds a Kalman filter inside the likelihood, is the accepted extension \cite{jategaonkar2006flight}.

\subsection{Variational System Identification}
Variational system identification occupies the space between deterministic OEM and filter-error methods. For a stochastic continuous-time model,
\[
d\mathbf{x}(t)=f(\mathbf{x}(t),\mathbf{u}(t);\theta_p)\,dt+G(\theta_p)\,dW(t),
\qquad
\mathbf{y}_k=h(\mathbf{x}_k,\mathbf{u}_k;\theta_p)+\mathbf{v}_k,
\]
the exact likelihood requires marginalizing over the latent state path $\mathbf{x}_{0:N}$. FEM approximates this marginalization using a Kalman-filter prediction model. Variational system identification instead introduces an assumed smoothing density $q(\mathbf{x}_{0:N})$ and maximizes an evidence lower bound,
\[
\mathcal{L}(\theta_p,q)
=
\mathbb{E}_{q}
\left[
\log p_{\theta_p}(\mathbf{x}_{0:N},\mathbf{y}_{0:N})
-\log q(\mathbf{x}_{0:N})
\right],
\]
over both the model parameters and the smoothing distribution. In Dutra's aircraft formulation \citep{dutra_variational_2025}, the assumed density is chosen to preserve a Markov structure so that the objective remains sparse and tractable. This makes VI a natural candidate for the mocap-output problem considered here: the measured outputs are $p_x,p_z,\theta$, the flight states are latent, and turbulence or unmodeled disturbances can be represented as process noise rather than being forced into measurement residuals.

The role of VI in this repository benchmark is therefore different from SINDy, PINNs, or UDEs. VI improves inference for a prescribed grey-box model in the presence of process and measurement noise; it does not automatically supply a new aerodynamic closure. If the nominal coefficient structure is incomplete, VI can produce better state smoothing and uncertainty estimates, but the missing aerodynamic physics must still be represented through additional parameters, basis functions, Gaussian-process terms, or neural residuals. A natural extension of the present benchmark is therefore a variational UDE or variational coefficient-closure model, where VI handles latent-state and process-noise uncertainty while the residual model handles structural aerodynamic error.

\subsubsection{Multiple--Shooting Formulation}

In the multiple--shooting implementation \cite{bock1984multiple,jategaonkar2006flight}, the time horizon is divided into $M$
shorter windows. The aerodynamic parameters and noise standard deviations remain
global, while each window has its own optimized shooting-node state. The decision
vector is
\[
    w_{\mathrm{MS}}
    =
    \begin{bmatrix}
        \boldsymbol{\theta}_p^\top &
        \mathbf{X}_0^\top &
        \mathbf{X}_1^\top &
        \cdots &
        \mathbf{X}_{M-1}^\top &
        \log\boldsymbol{\sigma}^\top
    \end{bmatrix}^{\!\top}.
\]
Within the $i$th window, the local trajectory is initialized from $\mathbf{X}_i$
and propagated as
\[
    \hat{\mathbf{x}}_{i,k+1}
    =
    \Phi(\hat{\mathbf{x}}_{i,k},\mathbf{u}_{i,k};
    \boldsymbol{\theta}_p,\Delta t),
    \qquad
    \hat{\mathbf{x}}_{i,0}=\mathbf{X}_i.
\]
Continuity between neighboring windows is enforced through equality constraints,
\[
    \mathbf{g}_i
    =
    \hat{\mathbf{x}}_{i,\mathrm{end}}
    -
    \mathbf{X}_{i+1}
    =
    \mathbf{0},
    \qquad i=0,\ldots,M-2.
\]
The corresponding MLE problem is
\[
\begin{aligned}
    \min_{\boldsymbol{\theta}_p,\,\{\mathbf{X}_i\},\,\log\boldsymbol{\sigma}}
    \quad
    J_{\mathrm{MS}}
    &=
    \sum_{j=1}^{4} N\log\sigma_j
    +
    \frac{1}{2}
    \sum_{i=0}^{M-1}
    \sum_{k=0}^{N_i-1}
    \sum_{j=1}^{4}
    \frac{
    \left(\hat{y}_{i,k,j}-y_{i,k,j}\right)^2
    }{\sigma_j^2} \\
    \text{s.t.}\quad
    \mathbf{g}_i &= \mathbf{0},
    \qquad i=0,\ldots,M-2.
\end{aligned}
\]
This introduces additional decision variables but allows the optimizer to leverage intermediate state information, which significantly improves robustness and convergence behavior. In the current benchmark, the generated comparison figures and tables are the authoritative output-error results; no separate hand-maintained SS/MS demonstration figure is included.
\subsection{SINDy and its Variations}
Sparse Identification of Nonlinear Dynamics (SINDy) searches for governing
equations by regressing measured state derivatives on a user-defined library of
candidate nonlinear functions. For the longitudinal benchmark, the dependent
variable for a fully black-box SINDy model would be the smoothed derivative matrix
\[
    \dot{X} =
    \begin{bmatrix}
    \dot{V} & \dot{\alpha} & \dot{\gamma} & \dot{Q}
    \end{bmatrix},
\]
and the candidate library is evaluated on
$[V,\alpha,\gamma,Q,T,\delta_e]$. The standard SINDy model is
\begin{equation}
    \dot{X} = \Theta(X,U)\,\Xi,
    \label{eq:sindy_standard}
\end{equation}
where each column of $\Xi$ contains the active terms for one state equation. A
LASSO-style sparse formulation can be written as
\begin{equation}
    \Xi^\star
    =
    \operatorname*{arg\,min}_{\Xi'}
    \left\|
    \Theta(X,U)\Xi' - \dot{X}
    \right\|_F^2
    +\lambda\|\Xi'\|_1 .
    \label{eq:sindy_lasso}
\end{equation}
The implementation used for the benchmark follows the common sequential
thresholded least-squares philosophy rather than a full convex LASSO solve, but
applies it to aerodynamic coefficients rather than directly to
\eqref{eq:sindy_standard}. The low-order aircraft terms are protected in the
regression, while additional nonlinear terms are sparse. This lets SINDy learn
the nominal-like coefficient parameters and the missing nonlinear aerodynamic
terms in a single fit. The same library is also used by a forward stepwise
symbolic-regression baseline, which remains a black-box derivative fit.

For aircraft data, the practical details are as important as the sparse
optimizer. Derivatives are obtained from smoothed mocap-derived state
histories; this is intentionally difficult
because motion-capture position and attitude must be differentiated to recover
airspeed, flight-path angle, angle of attack, and pitch rate. Hidden actuator
or SAFE-loop dynamics are also challenging because the candidate library sees
the external pilot command rather than the internal control signal. This is why
the benchmark uses structured coefficient SINDy: the known rigid-body balance is
built into the coefficient-inference step, low-order aircraft terms are retained
as protected regressors, and sparse discovery is reserved for missing nonlinear
aerodynamic terms.

\subsection{Gaussian-Process and Kernel Closure Methods}
Gaussian-process regression places a prior over the unknown residual
dynamics $f_{\mathrm{res}}(\mathbf{z})$ and conditions it on training
transitions, yielding both a posterior mean and a predictive variance
\cite{rasmussen2006gaussian}. In aircraft identification GP models have been
used as aerodynamic-coefficient surrogates and residual closures
\cite{hemakumara2013learning}, and GP state-space formulations with
variational inference extend this to latent states and process noise
\cite{frigola2014variational,doerr2018probabilistic}. The repository's
kernel rows implement the deterministic core of this idea: a radial-basis
kernel expansion over the heading- and position-invariant features
$\mathbf{z} = [u, v, w, \mathbf{g}_b, p, q, r, \mathbf{u}_{\mathrm{cmd}}]$,
\[
\hat f_{\mathrm{res}}(\mathbf{z}) = \sum_{i=1}^{N_c} w_i\,
\exp\!\Big(-\tfrac{1}{2}\,\big\Vert (\mathbf{z}-\mathbf{c}_i) \oslash
\boldsymbol{\ell} \big\Vert^2\Big),
\]
with $N_c$ centers subsampled from the training inputs, per-dimension length
scales $\boldsymbol{\ell}$ set from the feature standard deviations, and
ridge-regularized weights -- equivalent to the posterior mean of a sparse GP
with fixed hyperparameters. The residual corrects only the six dynamic
states around the nominal model; position and attitude integrate
kinematically. Two honest limitations relative to full GP practice are
stated in the method notes: hyperparameters are not optimized by marginal
likelihood, and the predictive variance is not propagated, so the rows
report point predictions only. Upgrading to variance-aware free runs (a
predictive tube rather than a line) is the natural extension for the
website's rollout display.

\subsection{Physics-Informed and Neural-Residual Methods}
For the longitudinal benchmark considered here, a direct inverse PINN that penalizes deviations from the same fixed aerodynamic coefficient model used by OEM is not the most appropriate use of physics information. The purpose of the benchmark is to evaluate identification methods when unmodeled aerodynamic forces, moments, actuator effects, or disturbances are present. If the PINN physics residual forces the trajectory to satisfy an incomplete nominal force and moment model, the residual acts as a regularizer toward the wrong dynamics and can suppress the missing effects the method is supposed to discover.

The useful distinction is between exact structure and aerodynamic closure assumptions. Kinematic identities and rigid-body balance laws should be enforced strongly or built into the parameterization. By contrast, aerodynamic force and moment functions should remain learnable because these are precisely the terms that may change outside the nominal model. A tractable PINN formulation for this benchmark therefore uses a neural trajectory $\widehat{\mathbf{x}}_\theta(t)$ and learned aerodynamic coefficient functions,
\[
\widehat{C}_{D,\phi}=\mathcal{N}_{D,\phi}(\widehat{\alpha},\widehat{q},\delta_e,\widehat{V}),\quad
\widehat{C}_{L,\phi}=\mathcal{N}_{L,\phi}(\widehat{\alpha},\widehat{q},\delta_e,\widehat{V}),\quad
\widehat{C}_{M,\phi}=\mathcal{N}_{M,\phi}(\widehat{\alpha},\widehat{q},\delta_e,\widehat{V}),
\]
rather than treating a low-order coefficient expansion as the true physics. The corresponding dimensional forces and moments are
\[
\widehat{D}=\bar{q}S\widehat{C}_{D,\phi},\qquad
\widehat{L}=\bar{q}S\widehat{C}_{L,\phi},\qquad
\widehat{M}=\bar{q}S\widehat{C}_{M,\phi},
\]
where $\bar{q}=\frac{1}{2}\rho \widehat{V}^2$. The differential-equation residual then enforces the longitudinal rigid-body equations,
\[
\begin{aligned}
r_V &= \dot{\widehat{V}}
 - \frac{1}{m}\left(-\widehat{D}+T\cos\widehat{\alpha}-mg\sin\widehat{\gamma}\right),\\
r_\gamma &= \dot{\widehat{\gamma}}
 - \frac{1}{m\widehat{V}}\left(\widehat{L}+T\sin\widehat{\alpha}-mg\cos\widehat{\gamma}\right),\\
r_q &= \dot{\widehat{q}}-\frac{\widehat{M}}{J_y}.
\end{aligned}
\]
The pitch kinematic relationship, such as $\dot{\alpha}=q-\dot{\gamma}$ for the standard longitudinal coordinates, is not an aerodynamic discovery target and should preferably be enforced by the state parameterization itself. If it is not hard-coded, it can be included as a high-weight residual, but this is a numerical convenience rather than a claim that the kinematics are uncertain physics.

The resulting training objective is
\[
\mathcal{L}(\theta,\phi)=
\frac{\lambda_{\mathrm{data}}}{N_y}\sum_i
\left\|\mathbf{H}\widehat{\mathbf{x}}_\theta(t_i)-\mathbf{y}_i\right\|_2^2
+\frac{\lambda_{\mathrm{dyn}}}{N_f}\sum_j
\left\|\mathbf{r}_{\mathrm{dyn}}(t_j)\right\|_2^2
+\lambda_{\mathrm{reg}}\mathcal{R}_{\mathrm{aero}}(\phi),
\]
where $\mathbf{r}_{\mathrm{dyn}}=[r_V,r_\gamma,r_q]^T$. The regularization term can encode weak aerodynamic priors such as smoothness, dynamic-pressure scaling, drag positivity, boundedness, and local derivative signs near trim. These priors are much weaker than enforcing a complete aerodynamic model and are therefore more consistent with the goal of identifying unmodeled force and moment behavior.

This setup places the present benchmark in the same family as the Non-Determinant and Modified Non-Determinant PINNs of Michek et al.~\citep{michek2026flight}, but at a deliberately smaller scale. Their work learns aerodynamic coefficients for a nonlinear 6-DOF simulation using many randomized flight intervals. In the legacy 3-DOF tier, the same conceptual role was isolated in a longitudinal problem with known nominal coefficients and hidden nonlinear residual terms, making it possible to compare PINN-style closure learning directly against OEM, SINDy, and UDE residual models.

A UDE formulation uses the same logic but augments a nominal coefficient model with learned residuals,
\[
\begin{aligned}
C_D &= C_{D,\mathrm{nom}}(x,u;\theta_p)+\Delta C_{D,\phi}(x,u),\\
C_L &= C_{L,\mathrm{nom}}(x,u;\theta_p)+\Delta C_{L,\phi}(x,u),\\
C_M &= C_{M,\mathrm{nom}}(x,u;\theta_p)+\Delta C_{M,\phi}(x,u).
\end{aligned}
\]
This is especially plausible for the present scope because the known model is used where it is trusted, while the neural residual captures model-form error. The identified parameters $\theta_p$ remain interpretable, and the residual terms provide a direct diagnostic for where the nominal aerodynamic model is inadequate.

This formulation is intentionally not a full aerodynamic PINN in the Navier-Stokes sense. Without geometry, flow-field states, boundary conditions, and a high-fidelity solver or data set, enforcing the Euler, RANS, or Navier-Stokes equations is not tractable in this repository benchmark. The PINN comparison is therefore limited to flight-dynamics-level grey-box learning, while high-fidelity CFD-constrained surrogate modeling is treated as future work.




\section{Method Selection Guidance}
The benchmark is designed to support method selection rather than identify a
single universally best algorithm. When direct measurements of
$V,\alpha,\gamma,Q$ are available and the primary mismatch is unknown
aerodynamics, coefficient-closure and residual-learning methods are attractive
because they preserve the rigid-body equations while learning the missing force
and moment terms. The supervised coefficient surrogate is the best possible
version of this idea in the synthetic benchmark, but it uses labels that real
flight tests do not provide; it should therefore be interpreted as an upper
bound. Flight-data-only UDE and PINN-style closure methods are more realistic
choices in this regime.

When only motion capture is available, the main difficulty changes. Algorithms
that require accurate state derivatives suffer from finite-difference noise,
smoothing lag, and loss of bandwidth. For this reason, mocap-output OEM and
variational or smoothing formulations are better
matched to the open-loop model-identification problem than methods that first reconstruct
$V,\alpha,\gamma,Q$ and then differentiate them. Earlier revisions also
reported an idealized direct-state diagnostic; the benchmark now evaluates only
the motion-capture observation policy expected in the PURT pylon-racing
measurement setup.

When SAFE is enabled, the estimator observes the external pilot command while
the plant responds to an internally reshaped elevator command. In this setting,
a purely aerodynamic residual model can fit trajectories while misattributing
hidden control action to aerodynamic effects. A hidden-input UDE or a
parameterized hidden-controller OEM is more conceptually appropriate. The OEM
hidden-controller row is deliberately optimistic because it assumes the analyst
has guessed the controller structure. Real proprietary controllers may contain
mode switches, saturations, rate limits, and recovery gates that are only
identifiable if the flight test excites those regions. Therefore, ordinary lap
data may be sufficient to learn a local predictive model but insufficient to
reverse-engineer the controller.

For frequency-domain identification, dedicated excitation matters. The compact
Frequency-Welch baseline is useful for checking bandwidth and coherence, but it
is not a substitute for a purpose-built CIFER-style campaign with carefully
designed multisines, sweeps, input separation, and coherence validation. The
sine-sweep datasets are included to make this distinction explicit.

The practical recommendation is therefore problem dependent. Use mocap-output
OEM or variational/filter-error methods when the measurement system is the
dominant limitation. Use UDE or PINN-style coefficient closures when the
measurement states are reliable but the aerodynamic model is incomplete. Use
GP/RBF coefficient closures when a kernel surrogate with smoother interpolation
and an eventual path to posterior uncertainty is preferred over a neural
closure. Use Koopman-EDMD, subspace, or other lifted/input-output surrogates
when predictive accuracy over the tested envelope is more important than
recovering aerodynamic parameters. Use local model stitching or LPV-style
scheduling when separate parts of the flight envelope are well represented by
simple models but one global model averages incompatible trim, stall-onset, and
recovery behavior. Use
hidden-input models when a proprietary autopilot reshapes commands. Use
frequency-domain tools when the experiment has been designed for spectral
identification. The value of the benchmark is that each method is exposed to
the same train/validation splits across these regimes, so these recommendations
can be revised quantitatively as additional methods and datasets are added.

The updated aggressive results also clarify the role of simple linear models.
Linear-SS is not invalid; it is an effective local predictor for small motions
and can remain competitive when mocap-derived state reconstruction dominates
the error budget. It should not, however, be interpreted as recovering the
nonlinear aerodynamics. When the maneuvers expose the stall and recovery
regimes, the nonlinear closure and residual-learning methods give substantially
lower validation error. This is the intended benchmark behavior: near-trim tests
verify that simple models are not unfairly penalized, while aggressive
initial-condition and maneuver variation reveal when global linearity is no
longer adequate.

\section{Repository Roadmap}
The immediate roadmap is repository-driven rather than submission-driven. The first priority is to keep the 6-DOF benchmark reproducible: generated datasets, figures, tables, website artifacts, and this reference document should be regenerated from the same command-line entry points. The second priority is to mature the remaining placeholder methods (the frequency-domain rows toward a CIFER-style identification, and the variational, PINN-closure, and neural-surrogate rows toward genuine implementations). A useful future extension is a separate F-16 model family based on the Stevens--Lewis table-lookup aerodynamics, which would test high-performance aircraft identification without replacing the present small-RC, mocap-motivated benchmark. The third priority is to stabilize the method-plugin interface so external contributors can add methods without editing the core comparison harness.

Future methods should be added only when the repository contains executable code, benchmark metadata, and validation results for the new row. Candidate additions include input-output nonlinear predictors, richer local/LPV model stitching, Bayesian uncertainty-aware variants of the existing grey-box methods, and high-fidelity aerodynamic surrogate closures. Until those implementations exist in the repository, this reference limits detailed method descriptions to the rows already represented by code and generated results.

\appendix
\section{Method Implementations and Results}
\label{app:implementation}

This appendix presents each implemented method as a self-contained unit ---
approach, implementation, and results --- so a reviewer can audit the code
against the stated algorithm and the claims against the generated numbers.
Unless stated otherwise, every score in this appendix is one of the
following precisely defined quantities, all computed on open-loop rollouts
that receive only an initial state and the recorded pilot commands (rule
R1). \emph{Full-state mean NRMSE} (the leaderboard ``validation score''):
for each of the 13 state channels, the root-mean-square error between the
rollout and the measured validation states, pooled over all validation
chunks and samples, is divided by that channel's min-to-max range in the
validation truth; the score is the unweighted mean of the 13 normalized
values. Quaternion signs are aligned to the truth before scoring (a gauge
fix, $q\equiv-q$); non-finite rollouts score $10^9$ and physically
diverged rollouts are flagged in a separate column rather than masked.
\emph{Pose RMSE} (reported in meters): the RMSE of the 3-D position
channels alone over the same rollouts, unnormalized. \emph{Grey-box
chunk NRMSE}: the same range-normalized construction applied to the
grey-box's 60\,Hz chunk rollouts over its twelve Euler states.
\emph{One-step dynamic residual}: the Euclidean norm of the model's
single-step prediction error restricted to the six dynamic states
(body velocities and rates, in their native mixed units), used only for
regime comparisons such as the saturation analysis. The five-second
free-run position metric used for the SAFE model is defined in its results
paragraph. ``Synthetic'' results average the four simulated scenarios and
``real'' results are the Sport Cub motion-capture datasets (4/17 and 5/22)
with mocap-derived states. The authoritative code is
\texttt{models/aircraft6dof/comparison\_suite.py},
\texttt{models/aircraft6dof/greybox\_oem\_fit.py},
\texttt{models/aircraft6dof/flight\_explorer\_export.py}, and
\texttt{site/src/explorer.js}.

\subsection{Common Experimental Rules}

Each rule was added in response to a verified failure mode and applies to all
methods unless stated.

\paragraph{R1 --- Validation data is initial conditions and inputs only.}
Every validation rollout receives one initial state and the pilot-command
history; no method filters, smooths, or conditions on validation
measurements mid-rollout, and measured states reappear only in scoring. The
one audit violation (the legacy 3-DOF suite defaulting to the post-controller
actuator signal $u_{\mathrm{act}}$) was fixed: $u_{\mathrm{cmd}}$ is the
default and $u_{\mathrm{act}}$ is documented as an oracle diagnostic.

\paragraph{R2 --- Known kinematics are never fitted.}
$\dot{\mathbf p}=R(q)\mathbf v$ and
$\dot q=\tfrac12 q\otimes\boldsymbol\omega$ are exact identities. Surrogates
fit only the six dynamic states $(\dot u,\dot v,\dot w,\dot p,\dot q,\dot
r)$; rollouts integrate position from the rotated body velocity and attitude
with an exact axis-angle quaternion step. Measured motivation: learned
position increments translated at 66\,m/s against a 12\,m/s velocity clamp
(impossible 134\,m errors in 3\,s), and linear features cannot represent the
bilinear $q\otimes\boldsymbol\omega$. Classical equation error likewise fits
only force and moment equations \cite{KleinMorelli2006}.

\paragraph{R3 --- Features are heading- and position-invariant.}
Surrogate features are $[u,v,w,\mathbf g_b,p,q,r,\mathbf u_{\mathrm{cmd}}]$
with $\mathbf g_b=R^{T}(q)\mathbf e_z$ the body-frame gravity direction:
flight dynamics are invariant to translation and heading, and attitude
enters only through gravity. Raw quaternion features let the 1.4k-sample
4/17 dataset overfit heading, and position features made rollouts
structurally divergent (worst published score 9{,}073 before the fix).

\paragraph{R4 --- Derivative and increment targets are smoothed.}
A Savitzky--Golay quadratic filter (50\,ms window, $\sim$12 samples) forms
derivative targets and smoothed-state differences form increment targets:
raw one-step differences of 240\,Hz mocap states have measured per-channel
signal-to-noise of only 0.24--0.77.

\paragraph{R5 --- Physical envelopes catch divergence honestly.}
6-DOF rollouts clamp speed and rates and keep clipped samples so the
diverged flag triggers honestly (in the legacy 3-DOF tier an exponential
blowup once posted an $\alpha$-RMSE of 1{,}160 radians as a finite,
unflagged score).

\paragraph{R6 --- Identical real-data preprocessing.}
Manual maneuvers split into quasi-steady lead-ins (pooled per-flight
data-compatibility trim-bias estimation) and 0.6--1.2\,s gap-free actuation
chunks with locally estimated initial states; train/validation is
round-robin per flight; mocap dropouts are never trained on or scored.
Result rows carry an implementation status (\emph{implemented},
\emph{analogue}, \emph{placeholder}) and a diverged flag.

\subsection{6DOF-NominalGreyBox}
\paragraph{Approach.} The repository's attached-flow quaternion rigid-body
model evaluated with nominal coefficients and the recorded pilot commands ---
no fitted parameters. On synthetic data the truth model is this nominal
grey-box plus hidden nonlinear residuals, so the row measures exactly the
unmodeled-physics gap every fitted method must close.
\paragraph{Implementation.} A parallel RK4 rollout of
\texttt{models/aircraft6dof/model.py} from the shared validation initial
states. The row appears only for synthetic scenarios: on real flights the
truth-minus-residual interpretation does not exist, so the row is omitted
(the model remains the internal base of the residual methods).
\paragraph{Results.} Synthetic 0.207 (direct and mocap) --- the floor that
separates ``learned something real'' (below) from ``fit noise'' (above).

\subsection{6DOF-GreyBoxOEM}
\paragraph{Approach.} Output-error estimation \cite{jategaonkar2006flight}
of a 22-coefficient physical model: lift, drag, and side-force expansions
($C_{L_0}, C_{L_\alpha}, C_{D_0}, C_{D_{CL^2}}, C_{Y_\beta}$, thrust gain
$K_T$) and roll/pitch/yaw moment-coefficient groups with $b/2V$ and
$\bar c/2V$ rate non-dimensionalization and surface-deflection terms. Each
training chunk is rolled out from its measured initial state and the
parameters jointly minimize position and Euler-attitude residuals weighted
by output sigmas --- segment-wise output error rather than multiple
shooting, since measured chunk initial conditions remove the need for
continuity variables.
\paragraph{Implementation.} CasADi RK4 at a 60\,Hz model rate inside scipy
trust-region least squares with physical bounds. Control-channel maps are
derived from name lists (a hardcoded map once silently swapped aileron and
elevator: scores barely moved, parameter signs were unphysical --- the
motivating incident for never hardcoding orders). Two physics bugs found by
review and fixed: thrust folded into the wind-axis rotation (a propeller is
body-fixed; worth 19\% of validation NRMSE) and bound starvation ---
the original $K_T$ cap of 0.85\,m/s$^2$ per unit throttle could not balance
cruise drag ($\approx$2.3 required), pinning $K_T$, $C_{L_0}$, and
$C_{L_\alpha}$; the fitter now records parameters within 1\% of the box.
$C_{D_{CL^2}}$ at its physical zero bound is accepted (induced drag is
weakly identifiable over the narrow $C_L$ range). A 25-parameter extension
adding servo/motor lag states and battery thrust decay was implemented and
gated out on held-out data (all three added parameters pinned at bounds,
validation worsened 0.106 $\rightarrow$ 0.130): the recorded 2--4\,rad/s
excitation cannot observe a $\sim$30\,ms servo pole, and maneuver
placement confounds the per-flight-time decay term.
\paragraph{Results.} Real data 0.158 (4/17) and 0.102 (5/22) --- the
reference physical model. Fitted values are physically plausible
($C_{L_0}=1.21$, $C_{L_\alpha}=3.59$ for a low-Reynolds cambered foam wing;
$K_T=1.36$; $C_{m_q}$-type damping $-24.7$) and the lateral control signs
match the real airframe (opposite the synthetic convention).
Cram\'er--Rao analysis identifies the lift line, damping/stiffness, and
control effectivenesses to 1--8\% relative standard deviation, finds
$C_{D_{CL^2}}$ unidentifiable (368\%), and exposes four strong couplings
($C_{L_0}\!\leftrightarrow\!C_{L_\alpha}$,
$K_{M_0}\!\leftrightarrow\!K_{M_\alpha}$,
$K_{L_0}\!\leftrightarrow\!K_{L_{\delta r}}$,
$K_{N_0}\!\leftrightarrow\!K_{N_{\delta r}}$) traced to the narrow
angle-of-attack band and the near-constant rudder trim. Not run on
synthetic scenarios (the spec encodes the Sport Cub airframe).

\subsection{6DOF-LinearSS}
\paragraph{Approach.} A global one-step affine predictor
$\mathbf x_{k+1}=W^{T}[\mathbf x_k,\mathbf u_k,1]$ in the prediction-error
tradition \cite{ljung1999system}: the strongest cheap baseline near an
operating envelope, with no claim of recovering aerodynamic structure.
\paragraph{Implementation.} Ridge regression on the training chunks;
open-loop validation with physical clamps. Positions remain in the feature
set by layout contract with the browser integrator; the near-identity
discrete map keeps them benign (unlike derivative/increment surrogates,
where they were removed under R3). The model inspector renders
$(W-I)/\Delta t$ as continuous-time equations.
\paragraph{Results.} Synthetic 0.376/0.388; real 0.207 (4/17) and 0.142
(5/22). On real laps it is hard to beat for short horizons but cannot
extrapolate maneuvers; the grey-box and UDE rows overtake it there.

\subsection{6DOF-RidgeResidual}
\paragraph{Approach.} The nominal grey-box RK4 step plus an affine one-step
residual --- the linear-residual version of the hybrid physics-plus-learning
idea, kept fast and deterministic.
\paragraph{Implementation.} Residual targets are next-state minus
nominal-step on the training chunks; ridge solve; rollout adds the residual
to the nominal step each sample.
\paragraph{Results.} Synthetic 0.260/0.280 (a solid improvement over the
0.207-floor nominal once its residual absorbs hidden stall terms ---
note the nominal row's 0.207 is the \emph{average} gap, while RidgeResidual
must also overcome mocap-derived initialization noise); real 0.693/0.247,
limited by its base model's wrong-signed lateral derivatives on the real
airframe --- the documented motivation for the UDE row's grey-box base.

\subsection{6DOF-EquationError-LS}
\paragraph{Approach.} Classical equation-error least squares
\cite{KleinMorelli2006,morelli_system_2002}: regress smoothed state
derivatives on instantaneous flight variables, fitting the force and moment
equations only --- fast, interpretable, and the canonical precursor that
seeds output-error methods.
\paragraph{Implementation.} Savitzky--Golay derivatives (R4) of the six
dynamic states (R2) on linear invariant features (R3) --- gravity-direction
components are what allow a \emph{linear} feature set to represent the
gravity term --- with a standardized ridge solve and explicit-Euler
validation inside the kinematic stepper.
\paragraph{Results.} Synthetic 3.49/3.76; real 1.16/2.23. The honest
structural ceiling of a single global \emph{linear} derivative model:
amplitude-dependent dynamics bias the fit, and integration accumulates the
bias over long rollouts. The quadratic-library SINDy row is the direct
upgrade and the score gap (2.2 vs 0.22 on 5/22) quantifies what the
quadratic terms buy.

\subsection{6DOF-SINDy}
\paragraph{Approach.} Sparse identification of nonlinear dynamics
\cite{brunton2016discovering,champion2019data}: sequentially thresholded
least squares selects a sparse model from a quadratic candidate library on
the invariant features, fitting smoothed derivatives of the six dynamic
states. The library is deliberately generic (structure \emph{discovery});
the aerodynamically structured alternative is the grey-box row.
\paragraph{Implementation.} \texttt{stlsq\_fit}: ridge fit
$\rightarrow$ per-output magnitude threshold (constant and linear blocks
protected) $\rightarrow$ ridge refit restricted to surviving features,
iterated to a fixed point. A single post-hoc prune is \emph{not} STLSQ
(it leaves dense-fit coefficient values in place) --- the backend label only
claims STLSQ since this was fixed. Typical sparsity: $\sim$210 of 630
coefficients.
\paragraph{Results.} Synthetic 0.296/0.298; real 0.126 (4/17) and 0.216
(5/22). On real data its discovered dominant terms are the quaternion-free
physical bilinears (e.g.\ $\dot q\leftarrow$ pitch-damping and elevator
terms), and it sits between the affine baseline and the physical grey-box
--- the expected position for unstructured discovery on limited data.

\subsection{6DOF-Symbolic-Stepwise}
\paragraph{Approach.} A sparse symbolic-regression-style baseline: the SINDy
library at milder sparsity, separating sparsity-level effects from library
effects. Labeled \emph{analogue}: it does not run statistical forward
selection (no PSE/PRESS criteria as in Klein's stepwise regression) or
genetic programming \cite{koza1990genetic}.
\paragraph{Implementation.} The same threshold-and-refit estimator at a
12\% retention quantile on smoothed increment targets.
\paragraph{Results.} Synthetic 0.243/0.240; real 0.602/0.648 --- denser
than SINDy and correspondingly more overfit on the small real datasets,
which is precisely the sparsity ablation the row exists to show.

\subsection{6DOF-Koopman-EDMD}
\paragraph{Approach.} An extended-DMD-style lifted predictor
\cite{koopman1931hamiltonian,williams2015data,schmid2010dmd}: a quadratic
dictionary approximates the action of the Koopman operator on state
observables. Labeled \emph{analogue}: lifted observables are not propagated
(predictions return to the state space each step), so the row tests the
EDMD feature class rather than the operator machinery.
\paragraph{Implementation.} Quadratic one-step increment predictor on
invariant features with heavy ridge (100$\times$), smoothed targets,
kinematic integration.
\paragraph{Results.} Synthetic 0.538/0.304; real 10.9 (4/17) and 1.0
(5/22). The 4/17 number is small-data quadratic overfit (1.4k samples vs
105 features) reported honestly; on 5/22 the same estimator is within
5$\times$ of the grey-box. Propagating a learned lifted state (true EDMD
with control) is the natural upgrade.

\subsection{6DOF-Subspace-Hankel}
\paragraph{Approach.} A short-memory linear predictor in the spirit of
subspace identification \cite{VanOverscheeDeMoor1996,verhaegen1992subspace,
HOKALMAN+1966+545+548}: a lag-3 history of the invariant state predicts the
next dynamic-state increment. Labeled \emph{analogue}: no Hankel SVD or
N4SID/MOESP realization is performed.
\paragraph{Implementation.} Ridge ARX on stacked invariant histories with
smoothed increment targets; the browser port re-seeds its lag window
whenever the incoming state is not its own last output (fresh anchor or
SAFE handoff).
\paragraph{Results.} Synthetic 0.240/0.239; real 0.110 and 0.111 ---
essentially tying the 22-parameter physical grey-box with a 32$\times$6
weight matrix. Two very different model classes agreeing is good evidence
the pipeline is honest, and that most of the predictable real-data signal
is short-memory linear; a true realization-based implementation is a
worthwhile follow-up.

\subsection{6DOF-GP-RBF}
\paragraph{Approach.} The deterministic core of a sparse Gaussian-process
residual closure (see the kernel formulation section):
$\hat f_{\mathrm{res}}(\mathbf z)=\sum_i w_i\exp(-\tfrac12\Vert(\mathbf z-
\mathbf c_i)\oslash\boldsymbol\ell\Vert^2)$ correcting the nominal model's
dynamic states \cite{rasmussen2006gaussian,hemakumara2013learning}.
\paragraph{Implementation.} 96 centers subsampled from training inputs,
per-dimension length scales from feature standard deviations, ridge
weights; equivalent to a sparse-GP posterior mean with fixed
hyperparameters. Stated limitations: no marginal-likelihood optimization,
no propagated predictive variance.
\paragraph{Results.} Synthetic 0.233/0.224 (the best synthetic fitted row:
smooth residual aerodynamics are exactly a kernel method's home turf); real
0.745/0.522, inheriting the nominal base-model mismatch on the real
airframe. Rebasing on the grey-box (as done for the UDE) and
variance-aware rollouts are the planned upgrades.

\subsection{6DOF-UDE-Residual}
\paragraph{Approach.} The universal-differential-equation premise: best
known physics plus a learned residual \cite{rackauckas2020universal}. On
real scenarios the base is the \emph{fitted grey-box}; on synthetic
scenarios, the nominal grey-box. Labeled \emph{analogue}: the residual is a
ridge map, not a trained network (the genuine ODE-in-the-loop torch UDE
is the separate 6DOF-UDE-NN row).
\paragraph{Implementation.} Residual targets are smoothed next dynamic
states minus the base one-step prediction; shrinkage scales with sample
count --- the standardized Gram grows with $n$, so the previous fixed
penalty was $\sim$$10^{-5}$ of the Gram diagonal, i.e.\ effectively
unregularized noise (the measured cause of this row's old 12.1 score). With
$\lambda=n$, a noise-dominated residual correctly shrinks toward zero.
Rollouts integrate kinematics from the \emph{corrected} dynamics (with the
base's uncorrected rates, attitude drift compounded --- a measured
failure).
\paragraph{Results.} Real 0.109 (4/17) and 0.095 (5/22) --- the best real
rows, edging the bare grey-box (0.158/0.102): the residual contributes a
small genuine correction once shrinkage prevents noise fitting. Synthetic
2.06/0.98 is poor because there the base is the nominal model and the
hidden stall residual is strongly state-dependent --- the honest message is
that residual quality tracks base quality.

\subsection{6DOF-Model-Stitching}
\paragraph{Approach.} A scheduled family of local affine models blended by
airdata (speed, $\alpha$, $\beta$), following stitched-model and LPV
practice \cite{Tischler1995SystemIM,toth2010lpv}.
\paragraph{Implementation.} Local centers from the training distribution;
per-center ridge fits; rollout blends predictions with scheduling weights.
\paragraph{Results.} Synthetic 0.357/0.378; real 3.57/0.39. The 4/17
failure is the scheduling variables leaving the training hull on held-out
maneuvers; trim-anchored stitching with derivative interpolation is the
literature upgrade.

\subsection{6DOF-UDE-NN}
\paragraph{Approach.} The genuine universal differential equation
\cite{rackauckas2020universal}: the fitted grey-box provides the known
physics as a differentiable torch port of its continuous dynamics, and a
two-layer tanh MLP on the invariant features adds a residual to the six
dynamic-state derivatives, integrated together through RK4.
\paragraph{Implementation.} Trained by simulation error over half-second
shooting segments of the measured training chunks (derivative targets are
never formed); the output layer initializes near zero so training starts
from the bare grey-box; validation is a frozen-network open-loop rollout
with the standard physical clamps; the horizon shrinks automatically for
short-chunk datasets; the row skips gracefully when torch is unavailable.
\paragraph{Results.} 0.104 on the 5/22 mocap data, improving its own base
through identical rollout machinery (0.115 $\rightarrow$ 0.109 in
isolation); the gap to the ridge-UDE row (0.095) is the torch rollout's
coarser 60\,Hz integration rather than the network.

\subsection{6DOF-EKF-ParamID (Filter Error)}
\paragraph{Approach.} The filter-error method
\cite{jategaonkar2006flight,kalman_new_1960}: an augmented-state extended
Kalman filter estimates the 22 grey-box parameters jointly with the twelve
Euler states over the manual training chunks.
\paragraph{Implementation.} Joseph-form updates; CasADi step and parameter
Jacobians; the parameter covariance carries across chunks while the state
covariance resets at each measured chunk start; the $\psi$ innovation is
wrapped; parameters clip to their physical bounds. No filtering touches
validation data: validation is a frozen-parameter open-loop rollout.
\paragraph{Results.} 0.098 on the 5/22 mocap data --- ahead of the
trust-region output-error fit (0.102) at roughly 1/30th the training cost
--- and the companion 6DOF-Fisher-UQ row carries the grey-box Cram\'er--Rao
uncertainty report (weak parameters and couplings in its notes).

\subsection{Placeholders}
Frequency-Welch and Frequency-Stitching (regularized realizations with no
spectral content), Variational-Mocap (smoothed weak-form ridge; 3.71/3.67
synthetic), PINN-Closure (a sparsified copy of the UDE weights), and
NN-Surrogate (a random-feature closed form) are labeled \emph{placeholder}
in every artifact, are excluded from headline claims, and exist to reserve
the rows. The former EKF-ParamID/Fisher-UQ/OEM-SS trio --- one ridge fit
reported under three names --- is gone: the first two are the real
implementations above on the real-flight scenarios, and OEM-SS is removed
outright since 6DOF-GreyBoxOEM is the output-error row. With the
filter-error method and the trained network now real, the remaining
roadmap placeholder is CIFER-style frequency-domain fitting.

\subsection{SAFE Closed-Loop Model}
\paragraph{Approach.} Stabilized (SAFE-mode) flight cannot be predicted by
a bare airframe: the hidden inner loop adds surface corrections absent from
the recorded sticks. The \emph{direct method} of closed-loop identification
\cite{ljung1999system} fits the closed loop itself --- and is shared by
every method, since composing the separately identified controller with
each airframe was measured worse (13.6\,m vs 2.9\,m at 5\,s; the
inverse-dynamics controller estimate is weakly conditioned).
\paragraph{Implementation.} A linear ridge fit on the heading- and
position-invariant state: $[u,v,w,\phi,\theta,p,q,r,\mathrm{stick},1]
\rightarrow$ next dynamics plus the heading increment $\Delta\psi$;
heading integrates the fitted increment and position integrates
$R(q)\mathbf v$ exactly (a raw-state affine fit cannot represent
heading-dependent position kinematics: 8.3\,m $\rightarrow$ 2.9\,m at 5\,s
from this change alone). After the one-step ridge fit, the same
13$\times$9 weights are re-optimized against one-second free-run shooting
residuals over the training windows (simulation-error minimization):
one-step regression on noisy states carries an attenuation bias that
compounds over a 1{,}200-step rollout, and the multi-step objective removes
it with no structure change. Stabilized spans are cut into ten-second
windows, round-robin per flight, every third held out; the autonomous
flight is excluded (its lateral commands bypass the recorded sticks --- the
rudder channel is constant while the aircraft banks).
\paragraph{Results.} The headline metric is defined as follows. The
tracked stabilized spans are cut into ten-second windows and every third
window per flight is held out of all fitting (11 held-out windows across
the piloted flights; 52k samples remain for training). For each held-out
window, the model is initialized from the measured state at the window
start and integrated forward \emph{open loop} for five seconds (1{,}200
steps at 240\,Hz), driven only by the recorded stick history --- no
measurement updates of any kind. The error for that window is the
Euclidean distance in meters between the predicted and mocap-measured 3-D
position at exactly $t=+5$\,s, and the reported number is the mean over
the 11 windows. By this metric the refined model scores 2.98\,m ---
roughly 13\% of the $\sim$23\,m the aircraft travels in five seconds at
lap speed --- versus 3.41\,m for the same structure fitted one-step only,
and 21.9\,m for a naive baseline that freezes the body velocity and turn
rate at the anchor and coasts along the resulting arc. Free runs fly
complete laps. Envelope-protection saturation was hypothesized to dominate
the residual error and tested: one-step residuals are 74\% higher near
command saturation, but only $\sim$4\% of stabilized samples are affected
and a Hammerstein extension (fitted command scale and envelope saturation
ahead of the linear loop) improved the held-out error by only 3\%, so the
linear structure is retained and the finding is recorded here. The
manufacturer's manual clarifies the mechanism: the SAFE unit (gyro and
accelerometer only, no airspeed sensor) limits pitch and roll angles and
self-levels on stick release (the switch has three positions ---
self-leveling, envelope-only, and manual --- of which the recordings use
only the two extremes; the recorded mode channel is strictly bimodal at
1000/2000\,$\mu$s and the converter now refuses sessions that use the
middle position; the pilot confirms the Panic Recovery trigger was never
used, so no unlabeled fourth controller regime exists in the recordings), so the pilot-perceived low-airspeed
unresponsiveness is not a control-law gate: the controller is
speed-blind, and the dependence is plant-side --- every commanded
deflection buys a moment scaling with dynamic pressure, and the
attitude-command law only ever requests the gain-limited deflection
$K_p\cdot(\mathrm{error})$ toward a clipped setpoint --- a small slice of
the available throw. At low speed that slice buys too little moment to
track even the clipped bank command within maneuver timescales, while
speed-independent propeller disturbance moments grow in relative terms.
The airframe itself retains ample authority: the pilot can bank steeply at
the same speeds in Experienced mode by commanding full deflection
directly, which is precisely what the gain-limited law never does. A
speed-gated aileron-authority feature was tested accordingly and rejected
(held-out error worsened, 2.98\,m $\rightarrow$ 3.03--3.11\,m), and the
measured stick-to-roll coupling collapse below $\approx$4.5\,m/s is
explained by the same mechanism. The
fitted weights read as physics in the model inspector:
$\dot\phi\approx1.00\,p$, $\dot u\approx14.2-11.5\,\theta$ (thrust trim and
the $-g\theta$ gravity coupling), and $\dot p\approx-25\,\phi$ exposing the
SAFE self-leveling stiffness.

\subsection{SAFE Controller Grey-Box}
\paragraph{Approach.} A guessed attitude-command structure with fitted
gains and saturations: $\delta=K_p(\mathrm{sat}(c\cdot\mathrm{stick},
\pm\ell)-\mathrm{att})-K_d\cdot\mathrm{rate}+b$ per axis (rudder: stick
gain, yaw-rate gain, bias).
\paragraph{Implementation.} Staged grid search over $(c,\ell)$ with Huber
regression for the gains, fitted from stabilized segments by inverse
dynamics against the nominal airframe. The law is
proportional--derivative only --- no integral action --- consistent with
three observations: the transmitter carries manual elevator trim buttons
(integral action would make them redundant in stabilized mode), the fitted
closed-loop residuals show no unexplained slow mode, and attitude
integrators wind up in exactly the low-authority regime this airframe
occupies. A ten-second ground test discriminates directly: hold the
airframe tilted --- a ramping surface deflection indicates integral action,
a steady deflection pure PD.
\paragraph{Results.} Retained for interpretability and as the free-run
fallback; identified against the nominal airframe, which is why its
composition with each method's airframe currently loses to the direct
closed-loop fit --- re-identification against the fitted grey-box is the
planned route to per-method SAFE handoff.

\subsection{Legacy 3-DOF Implementation Decisions}
The legacy 3-DOF formulations appear in the supplemental sections; decisions
worth review: the suite OEM is CasADi multiple shooting (soft continuity
residuals, IPOPT budget 200, convergence status recorded; non-converged
solves are not published); the filter-error EKF (augmented state, Joseph
form, CasADi Jacobians) trains on training data only with frozen-$\theta$
open-loop validation; the torch UDE and PINN train RK4-in-the-loop over
multiple-shooting segments so derivative targets are never formed (the
benchmark sample period is 0.02\,s for Nyquist and RK4-stability reasons);
the SINDy library replaces trigonometric features with Taylor residuals to
break collinearity; trim solves the full equilibrium including lift
(residual $\sim$$10^{-16}$); the ETFE uses the correct
$\mathrm{csd}(\mathrm{input},\mathrm{output})$ ordering with input-power
weighting not claimed to be coherence, and unstable identified poles are
reflected into the left half-plane before open-loop simulation; and R5's
envelope applies.

\subsection{Browser Ports}
Every exported model has a JavaScript stepper implementing the same
equations as its Python counterpart; the grey-box port is verified against
the CasADi RK4 step to machine precision and all steppers are smoke-tested
for bounded, finite free runs from real anchors. Steppers share the exact
kinematic step (R2), the invariant feature map (R3), and the physical
clamps (R5); the free-run loop hands off between the manual-segment
airframe and the SAFE closed-loop model on the recorded mode channel and
stops at ground contact.

\section{Benchmark Result Tables}
\label{app:benchmark_tables}
The main text uses figures to summarize trends. The following tables retain the
complete sorted numerical results used to generate those figures. Direct-state
tables use simulated noisy $V,\alpha,\gamma,Q$ measurements. Mocap-derived
tables use states reconstructed from 100 Hz position and pitch-attitude
measurements. Lower validation score is better.

\subsection{6-DOF Baseline Results}
\IfFileExists{generated/aircraft6dof_method_comparison_table.tex}{%
\input{generated/aircraft6dof_method_comparison_table}
}{}

\IfFileExists{generated/shared_method_comparison_direct_table.tex}{%
\input{generated/shared_method_comparison_direct_table}
}{}
\IfFileExists{generated/shared_method_comparison_mocap_table.tex}{%
\input{generated/shared_method_comparison_mocap_table}
}{}

\end{document}
